\documentclass[11pt]{article}

\usepackage[a4paper,margin=1in]{geometry}
\usepackage{amsmath,amssymb,amsfonts,amsthm}
\usepackage{physics}
\usepackage{bm}
\usepackage{hyperref}
\usepackage{graphicx}
\usepackage{cite}

\title{
Prime-Structured Tunneling in a Multiplicity Bundle:\\
Connecting Four-Prime Newton Holography to Tabletop Experiments
}

\author{
  % Your name here
}

\date{\today}

\begin{document}

\maketitle

\begin{abstract}
We introduce a prime-labeled multiplicity field \(\Psi(x,t)\) whose
prime-structured potential
\(V_{\rm prime}(x) = \sum_k f_k(x) p_k^{-\sigma}\) deforms both quantum
tunneling barriers and the effective gravitational action. In a
four-prime truncation compatible with the four-prime Newton holography
model, the latent fourth prime enforces holographic embedding stability.
Within a one-dimensional double-well setting, WKB tunneling
probabilities acquire arithmetic modulations in energy, while the
variation of the tunneling probability with respect to curvature and
topological data induces an effective response \(\Delta\chi_{\rm eff}\)
in topology-sensitive invariants. The construction is directly
implementable in the AZ--TFTC one-dimensional Hamiltonian, yielding
falsifiable predictions for cavity spectra, Casimir deviations, and
curvature-weighted frequency shifts. This prime-structured tunneling
module thereby unifies Newton holography, Zeta--Schr\"odinger Dynamics,
and topological quantum effects within the Multiplicity Computing
Paradigm.
\end{abstract}

\tableofcontents

\section*{Executive Summary}

Multiplicity Theory treats physical and informational systems as
prime-structured multiplicity spaces rather than as collections of
independent point-like entities.\cite{hund_tunneling,multiplicity_overview}
In this report we formalize \emph{prime-structured tunneling} as the
first self-contained, simulation-ready module that connects this
philosophy to experimentally accessible quantum tunneling phenomena.

The core object is a prime-labeled multiplicity field
\(\Psi(x,t) = \sum_k c_k(x,t)\ket{p_k}_x\) defined on a one-dimensional
reduction of spacetime, with each fiber carrying a separable Hilbert
space basis labeled by the ordered primes.\cite{tunneling_topology_draft}
A prime-structured potential
\(V_{\rm prime}(x) = \sum_k f_k(x)p_k^{-\sigma}\) deforms a standard
double-well background \(V_0(x)\), producing an effective barrier
\(V(x)=V_0(x)+V_{\rm prime}(x)\) that directly enters the WKB tunneling
exponent.

In a four-prime truncation compatible with the four-prime Newton
holography model, the fourth latent prime stabilizes the embedding via a
simple closure condition on the prime weights. This yields arithmetic,
log-periodic structure in the tunneling probability \(T(E)\) and in the
associated spectral data of the AZ--TFTC Hamiltonian. Differentiating
the WKB exponent with respect to curvature-like data provides a concrete
expression for an effective topological response
\(\Delta\chi_{\rm eff}\), upgrading earlier heuristic links between
tunneling and topological invariants.

The module is directly simulatable in existing AZ--TFTC and
Zeta--Schr\"odinger Dynamics code, and generates falsifiable tabletop
predictions: prime-locked sidebands in cavity spectra, log-periodic
ripples in Casimir forces, and curvature-weighted frequency shifts in
tunneling-dominated regimes.

\section{Prime-Labeled Multiplicity Field}

\subsection{Multiplicity bundle and fibers}

Let \(M\) denote a (possibly reduced) spacetime manifold, and consider a
complex vector bundle \(\pi:\mathcal{E} \to M\) whose fibers
\(\mathcal{E}_x\) are separable Hilbert spaces equipped with an
orthonormal basis labeled by the ordered prime numbers.\cite{tunneling_topology_draft}
For each \(x\in M\),
\begin{equation}
  \mathcal{E}_x \cong \ell^2(\mathbb{P}) =
  \left\{ (c_k)_{k\ge 1} : \sum_{k=1}^\infty |c_k|^2 < \infty \right\},
\end{equation}
with canonical basis \(\{\ket{p_k}_x\}_{k\ge 1}\), where
\(\mathbb{P} = \{p_1,p_2,p_3,\dots\} = \{2,3,5,\dots\}\).

\begin{definition}[Prime-labeled multiplicity field]
A \emph{multiplicity field} on \(M\) is a section
\(\Psi:M\times \mathbb{R}\to \mathcal{E}\) of the form
\begin{equation}
  \Psi(x,t) = \sum_{k=1}^\infty c_k(x,t)\, \ket{p_k}_x,
  \label{eq:psi_def}
\end{equation}
such that for each \((x,t)\),
\begin{equation}
  \lim_{N\to\infty} \sum_{k=1}^{N} |c_k(x,t)|^2 < \infty.
  \tag{\(\star\)}
\end{equation}
We refer to \((\star)\) as the \emph{prime-sector square-summability}
constraint or \emph{lawfulness constraint}.
\end{definition}

The lawfulness constraint ensures that \(\Psi(x,t)\) lies in the Hilbert
fiber \(\mathcal{E}_x\) and that associated operators remain
Hilbert--Schmidt for suitable decay exponents.\cite{tunneling_topology_draft}

\subsection{Four-prime truncation and Newton holography}

For the purposes of connecting to the four-prime Newton holography
model, we consider a finite-prime truncation
\(\mathcal{P}_4 = \{2,3,5,7\}\). In this setting, the multiplicity field
reduces to
\begin{equation}
  \Psi^{(4)}(x,t) = \sum_{k=1}^{4} c_k(x,t)\, \ket{p_k}_x,
  \qquad p_1=2,\ p_2=3,\ p_3=5,\ p_4=7.
\end{equation}

The three smallest primes play the role of visible degrees of freedom,
while the fourth (latent) prime \(p_4=7\) encodes a hidden holographic
dimension that stabilizes the embedding.\cite{newton_holography_draft}
This stability is captured by a simple normalization condition on the
prime weights,
\begin{equation}
  \sum_{k=1}^{4} w_k\,p_k^{-\sigma} = 1,
  \qquad \sigma > \tfrac{1}{2},
  \label{eq:closure_condition}
\end{equation}
which can be derived from the requirement that an associated kernel
operator \(K\) be Hilbert--Schmidt in the four-prime sector.\cite{newton_holography_draft}

\section{Prime-Structured Tunneling Potential}

\subsection{Background double-well potential}

We consider a one-dimensional reduction of the system, either as a
tunneling coordinate in a more complex configuration space or as the
effective coordinate of an AZ--TFTC cavity mode.\cite{tunneling_topology_draft}
The background potential is chosen as a symmetric double well,
\begin{equation}
  V_0(x) = V_b\left(1 - \frac{x^2}{a^2}\right)^2 - E_0,
  \label{eq:double_well}
\end{equation}
with fixed parameters
\begin{equation}
  V_b = 1,\quad a = 2,\quad E_0 = 0.3.
\end{equation}
This potential has minima near \(\pm a\) and a central barrier of
height approximately \(0.7\) above the reference, so energies
\(E\in[0,0.6]\) lie in a tunneling-dominated regime.

\subsection{Prime-structured correction}

We augment \(V_0\) by a prime-structured correction
\(V_{\rm prime}(x)\) of the form
\begin{equation}
  V_{\rm prime}(x) = \sum_{k=1}^\infty f_k(x)\,p_k^{-\sigma},
  \qquad \sigma>0,
  \label{eq:Vprime_general}
\end{equation}
where \(f_k(x)\) encodes local geometric and multiplicity data.\cite{tunneling_topology_draft}
In the four-prime truncation compatible with Newton holography, we
restrict to
\begin{equation}
  V_{\rm prime}^{(4)}(x) = \sum_{k=1}^{4} f_k(x)\,p_k^{-\sigma},
  \qquad p_1=2,\ p_2=3,\ p_3=5,\ p_4=7.
  \label{eq:Vprime_4}
\end{equation}

We choose a concrete, curvature-coupled and oscillatory ansatz:
\begin{align}
  f_k(x) &= A_k\,R(x) + B_k \cos(\omega_k x + \phi_k),
  \label{eq:fk_def} \\
  R(x) &\approx x^2,
\end{align}
where \(R(x)\) plays the role of a one-dimensional toy analog of scalar
curvature.\cite{tunneling_topology_draft}

The coefficients and frequencies are fixed as
\begin{align}
  &A = [0.15, 0.12, 0.08, 0.05],\quad
    B = [0.08, 0.06, 0.04, 0.03],\nonumber\\
  &\omega_k = \frac{2\pi}{\log p_k},\quad
    \phi_k = 0,\quad \sigma = 0.8.
\end{align}
The decay exponent \(\sigma=0.8\) satisfies \(\sigma>1/2\), ensuring
Hilbert--Schmidt convergence of associated prime-weighted kernel
operators.\cite{tunneling_topology_draft}

\subsection{Total effective potential}

The total potential is then
\begin{equation}
  V(x) = V_0(x) + V_{\rm prime}^{(4)}(x),
  \label{eq:V_total}
\end{equation}
with \(V_0\) given by \eqref{eq:double_well} and
\(V_{\rm prime}^{(4)}\) by \eqref{eq:Vprime_4}--\eqref{eq:fk_def}. This
potential enters both the Schr\"odinger operator used in AZ--TFTC
simulations and the WKB analysis of tunneling.

\section{WKB Tunneling Probability and Prime Dependence}

\subsection{Standard WKB formula}

Consider a non-relativistic particle of mass \(m\) in one dimension,
governed by
\begin{equation}
  -\frac{\hbar^2}{2m}\frac{d^2\psi}{dx^2} + V(x)\psi = E\psi,
\end{equation}
where \(V(x)\) is given by \eqref{eq:V_total}. For an energy
\(E < V_{\text{max}}\) below the barrier, the WKB approximation yields
the tunneling probability
\begin{equation}
  T(E) \approx
  \exp\left[
    -\frac{2}{\hbar}
    \int_{x_L(E)}^{x_R(E)}
    \sqrt{2m\left(V(x) - E\right)}\,dx
  \right],
  \label{eq:T_WKB}
\end{equation}
where \(x_L(E)\) and \(x_R(E)\) are the classical turning points
defined by \(V(x)=E\).

\subsection{Prime-structured WKB exponent}

Substituting the explicit prime-structured potential \eqref{eq:V_total}
into \eqref{eq:T_WKB}, we obtain
\begin{equation}
  T(E) \approx
  \exp\left[
    -\frac{2}{\hbar}
    \int_{x_L(E)}^{x_R(E)}
    \sqrt{
      2m\left(
        V_0(x) + \sum_{k=1}^{4} f_k(x)p_k^{-\sigma} - E
      \right)
    }\,dx
  \right].
  \label{eq:T_prime_structured}
\end{equation}

To make the dependence on the prime weights explicit, define
\begin{equation}
  W(x) = V_0(x) - E, \qquad
  G(x) = \sum_{k=1}^4 f_k(x)p_k^{-\sigma},
\end{equation}
so that \(V(x)-E = W(x) + G(x)\). For sufficiently small or moderate
prime corrections (relative to the barrier height), we can expand the
square root to first order in \(G\),
\begin{align}
  \sqrt{W(x)+G(x)} 
  &\approx \sqrt{W(x)} +
    \frac{1}{2}\frac{G(x)}{\sqrt{W(x)}} + \mathcal{O}(G^2).
\end{align}
This yields
\begin{align}
  \ln T(E)
  &= -\frac{2}{\hbar}\int_{x_L}^{x_R}
     \sqrt{2m\left(W(x)+G(x)\right)}\,dx \nonumber\\
  &\approx -\frac{2}{\hbar}\int_{x_L}^{x_R}
     \sqrt{2m W(x)}\,dx
     -\frac{2}{\hbar}\int_{x_L}^{x_R}
     \frac{1}{2}\sqrt{\frac{2m}{W(x)}}\,G(x)\,dx
     + \mathcal{O}(G^2)\nonumber\\
  &= \ln T_0(E) +
     \delta \ln T(E) + \mathcal{O}(G^2),
\end{align}
where
\begin{align}
  \ln T_0(E)
  &:= -\frac{2}{\hbar}\int_{x_L}^{x_R}
      \sqrt{2m W(x)}\,dx,
  \\
  \delta \ln T(E)
  &:= -\frac{1}{\hbar}
      \int_{x_L}^{x_R}
      \sqrt{\frac{2m}{W(x)}}\,
      \sum_{k=1}^{4} f_k(x)p_k^{-\sigma}\,dx.
\end{align}
Thus, to leading order, prime structure modifies \(\ln T\) linearly via
the \(f_k(x)p_k^{-\sigma}\) contributions.

\subsection{Sensitivity to prime-sector coefficients}

We can compute the functional derivatives of \(\ln T\) with respect to
the geometric coefficients \(f_k\). Formally,
\begin{equation}
  \frac{\partial \ln T}{\partial f_k}
  \approx
  -\frac{1}{\hbar}
  \int_{x_L}^{x_R}
  \sqrt{\frac{2m}{W(x)}}\,p_k^{-\sigma}\,dx,
  \label{eq:dlnT_dfk}
\end{equation}
so that
\begin{equation}
  \frac{\partial T}{\partial f_k}
  = T\,\frac{\partial \ln T}{\partial f_k}
  \approx
  -\frac{T}{\hbar}
  \int_{x_L}^{x_R}
  \sqrt{\frac{2m}{W(x)}}\,p_k^{-\sigma}\,dx.
  \label{eq:dT_dfk}
\end{equation}
These expressions can be evaluated numerically for each \(k\) and each
energy \(E\), providing a direct measure of how variations in the
prime-coupled geometry affect tunneling probabilities.

\section{Topological Response and Curvature Coupling}

\subsection{Prime-modified Einstein equation (toy model)}

In the full Multiplicity framework, prime-number oscillations contribute
an additional term \(P_{\mu\nu}\) to the Einstein field equations, giving
\begin{equation}
  R_{\mu\nu} - \frac{1}{2}g_{\mu\nu}R + \Lambda g_{\mu\nu}
  + P_{\mu\nu}
  = \frac{8\pi G}{c^4} T_{\mu\nu},
  \label{eq:einstein_prime}
\end{equation}
where \(P_{\mu\nu}\) is modeled as a prime-weighted oscillatory
correction arising from the multiplicity field.\cite{tunneling_topology_draft}
In the one-dimensional reduction, we capture its effect via the
curvature-like scalar \(R(x)\) appearing in \(f_k(x)\).

\subsection{Effective topological response \(\Delta\chi_{\rm eff}\)}

Topological invariants such as the Euler characteristic \(\chi(M)\) are
functional of the curvature and global structure of the manifold. In
this prime-structured tunneling module, we introduce an \emph{effective}
topological response \(\Delta\chi_{\rm eff}\) defined by the sensitivity
of tunneling to curvature-like data.\cite{tunneling_topology_draft}

Let \(\mathcal{I}[g]\) denote a scalar constructed from the metric,
e.g.\ a scalar curvature density or Euler density. Assuming that the
geometric coefficients \(f_k(x)\) depend on \(\mathcal{I}[g]\),
we write
\begin{equation}
  \Delta \chi_{\rm eff}
  \approx \sum_{k=1}^{4}
  \left(
    \frac{\partial T}{\partial f_k}
  \right)
  \left(
    \frac{\partial f_k}{\partial \mathcal{I}}
  \right),
  \label{eq:Delta_chi_eff}
\end{equation}
where \(\partial T/\partial f_k\) is given by \eqref{eq:dT_dfk}. For the
explicit ansatz \(f_k(x)=A_k R(x) + B_k \cos(\omega_k x)\) and
\(\mathcal{I} \equiv R(x)\), we have
\begin{equation}
  \frac{\partial f_k}{\partial \mathcal{I}} = A_k.
\end{equation}
In this simplified setting, \eqref{eq:Delta_chi_eff} becomes
\begin{equation}
  \Delta \chi_{\rm eff}
  \approx \sum_{k=1}^{4}
  A_k\,
  \frac{\partial T}{\partial f_k},
\end{equation}
which can be computed numerically for each energy \(E\) and parameter
set, thereby turning the earlier heuristic relationship between
tunneling probability and topological change into a computable response
functional.

\section{AZ--TFTC Implementation and Tabletop Signatures}

\subsection{Discretized Hamiltonian with prime-structured potential}

The AZ--TFTC framework employs a one-dimensional Hamiltonian, which in a
finite-difference representation takes the form
\begin{equation}
  H = -\frac{\hbar^2}{2m} D^{(2)} + V_{\rm grid},
\end{equation}
where \(D^{(2)}\) is the discretized second-derivative operator and
\(V_{\rm grid}\) is the potential evaluated on a spatial grid
\(x \in [x_{\min}, x_{\max}]\).\cite{AZ_TFTC_draft}

Given grid points \(x_j\), the prime-structured potential is implemented
as
\begin{equation}
  V_{\rm grid}(x_j) = V_0(x_j) + V_{\rm prime}^{(4)}(x_j),
\end{equation}
with \(V_0\) and \(V_{\rm prime}^{(4)}\) as in
\eqref{eq:double_well}--\eqref{eq:Vprime_4}.

\subsection{Numerical pseudocode}

A minimal Python/NumPy implementation of the potential is:
\begin{verbatim}
import numpy as np

# Grid and constants (example)
x = np.linspace(-5.0, 5.0, 2000)
Vb, a, E0 = 1.0, 2.0, 0.3

# Background double-well
V0 = Vb * (1 - x**2 / a**2)**2 - E0

# Prime-structured correction
Vprime = np.zeros_like(x)
primes = [2, 3, 5, 7]
A = [0.15, 0.12, 0.08, 0.05]
B = [0.08, 0.06, 0.04, 0.03]
sigma = 0.8

for i, p in enumerate(primes):
    omega = 2.0 * np.pi / np.log(p)   # ω_k = 2π / log p_k
    fk = A[i] * (x**2) + B[i] * np.cos(omega * x)  # φ_k = 0
    Vprime += fk / (p**sigma)

V_total = V0 + Vprime
# V_total now goes into the finite-difference / MPS Hamiltonian
\end{verbatim}

This pseudocode can be inserted directly into the AZ--TFTC notebook to
generate spectra, tunneling amplitudes, and Casimir-like observables in
the presence of prime-structured corrections.

\subsection{Predicted signatures}

Within this setup, we expect:
\begin{itemize}
  \item \textbf{Prime-locked sidebands in cavity spectra}, arising from
  the log-prime frequencies \(\omega_k = 2\pi/\log p_k\) entering the
  oscillatory part of \(f_k(x)\).
  \item \textbf{Log-periodic ripples in tunneling rates and Casimir
  forces}, as a function of energy or separation, due to the
  \(p_k^{-\sigma}\) weighting and the interference between prime
  contributions.
  \item \textbf{Curvature-weighted frequency shifts}, detectable when
  the curvature-like term \(R(x)\) is varied (e.g.\ by modifying the
  effective geometry), which can be traced via the response
  \(\Delta\chi_{\rm eff}\).
\end{itemize}
These provide concrete, falsifiable predictions that connect the
Multiplicity framework to tabletop experiments in quantum tunneling and
cavity QED.\cite{tunneling_topology_draft,AZ_TFTC_draft,zeta_schrodinger_draft}

\section{Discussion and Outlook}

We have constructed a self-contained, prime-structured tunneling module
based on a prime-labeled multiplicity field, a four-prime truncated
potential, and a WKB analysis sensitive to curvature and topological
data. This module:
\begin{itemize}
  \item Embeds naturally into the four-prime Newton holography model,
  with the latent fourth prime stabilizing the embedding via
  \eqref{eq:closure_condition}.
  \item Implements directly in AZ--TFTC and Zeta--Schr\"odinger Dynamics
  as a modified potential term, enabling numerical exploration with
  modest computational resources.
  \item Provides a concrete formula for an effective topological
  response \(\Delta\chi_{\rm eff}\) driven by prime-weighted tunneling.
\end{itemize}

Future work includes:
\begin{itemize}
  \item Extending from the 1D toy curvature \(R(x)\approx x^2\) to
  higher-dimensional and more realistic geometric backgrounds.
  \item Investigating the relation between prime-structured tunneling
  and genuine topological invariants (Euler, Pontryagin) in full
  four-dimensional spacetimes.
  \item Coupling this tunneling module to time-crystal phases of the
  multiplicity field, thereby unifying non-equilibrium temporal order
  and prime-structured barriers.
\end{itemize}

\appendix

\section{Explicit WKB Derivation for Prime-Structured Potential}
\label{app:wkb}

Starting from the Schr\"odinger equation with potential
\(V(x)=V_0(x)+\sum_k f_k(x)p_k^{-\sigma}\), we define the action
\begin{equation}
  S(x) = \int^x \sqrt{2m\left(V(y)-E\right)}\,dy,
\end{equation}
and write the WKB wavefunction in the classically forbidden region
\(x\in[x_L,x_R]\) as
\begin{equation}
  \psi(x) \approx \frac{C}{\sqrt[4]{2m\left(V(x)-E\right)}}\,
  \exp\left(-\frac{1}{\hbar}S(x)\right).
\end{equation}
Matching across the turning points yields the standard result
\eqref{eq:T_WKB}. The explicit dependence on \(f_k\) then follows by
differentiating \(S(x)\) with respect to \(f_k\) under the integral, as
outlined in Section~4.

\section{Numerical Implementation Details}
\label{app:numerics}

This appendix can include:
\begin{itemize}
  \item Details of the finite-difference discretization of the kinetic
  term and boundary conditions.
  \item Pseudocode or actual code snippets for computing eigenvalues,
  tunneling probabilities, and response functions.
  \item Example plots of \(V(x)\), \(|\psi(x)|^2\), and
  \(T(E)\) for representative parameter choices.
\end{itemize}

\appendix
\section*{Mathematical Appendix}
\addcontentsline{toc}{section}{Mathematical Appendix}

In this appendix we collect rigorous statements and proofs for the
operator-theoretic aspects of the prime-structured tunneling module.
We work in the one-dimensional setting described in the main text, with
potential
\[
  V(x) = V_0(x) + V_{\rm prime}^{(4)}(x),
\]
where \(V_0\) is the double-well potential \eqref{eq:double_well} and
\(V_{\rm prime}^{(4)}\) is the four-prime correction
\eqref{eq:Vprime_4}--\eqref{eq:fk_def}.

Throughout we assume units where \(\hbar = 1\) unless otherwise noted.

\section{Hilbert Space and Prime-Labeled Multiplicity Sector}

\subsection{Separable prime-labeled fiber}

Let \(\mathbb{P} = \{p_1,p_2,p_3,\dots\} = \{2,3,5,\dots\}\) denote the
ordered set of prime numbers. We define the Hilbert space
\begin{equation}
  \mathcal{H}_{\rm prime} \coloneqq \ell^2(\mathbb{P})
  = \left\{ (c_k)_{k\ge 1} : \sum_{k=1}^{\infty} |c_k|^2 < \infty \right\},
\end{equation}
equipped with the standard inner product
\begin{equation}
  \braket{c}{d}
  = \sum_{k=1}^{\infty} \overline{c_k}\,d_k.
\end{equation}
We denote by \(\{\ket{p_k}\}_{k\ge 1}\) the canonical orthonormal basis
defined via
\begin{equation}
  \ket{p_k} \longleftrightarrow (0,\dots,0,1,0,\dots),
\end{equation}
with a \(1\) in the \(k\)-th position and zeros elsewhere. Then
\(\mathcal{H}_{\rm prime}\) is separable, and every vector
\(\psi \in \mathcal{H}_{\rm prime}\) has the expansion
\(\psi = \sum_{k=1}^\infty c_k \ket{p_k}\) with
\(\sum_k|c_k|^2 < \infty\).

In the main text, for each spatial point \(x\) the fiber
\(\mathcal{E}_x\) is taken to be isomorphic to \(\mathcal{H}_{\rm prime}\),
and the multiplicity field is a section
\(\Psi(x,t) = \sum_k c_k(x,t)\ket{p_k}_x\) satisfying the prime-sector
square-summability constraint.

\subsection{Four-prime truncation}

For the four-prime Newton holography truncation we restrict attention to
the finite-dimensional subspace
\begin{equation}
  \mathcal{H}_{\rm prime}^{(4)}
  \coloneqq \operatorname{span}\{\ket{p_1},\ket{p_2},\ket{p_3},\ket{p_4}\}
  \cong \mathbb{C}^4,
\end{equation}
with \(p_1=2,p_2=3,p_3=5,p_4=7\).

Any vector \(\psi^{(4)} \in \mathcal{H}_{\rm prime}^{(4)}\) can be
written as
\begin{equation}
  \psi^{(4)} = \sum_{k=1}^4 c_k \ket{p_k},
  \qquad c_k\in\mathbb{C},
\end{equation}
and the lawfulness constraint is automatically satisfied.

\section{Prime-Weighted Multiplication Operators and Norm Bounds}

We now formalize the role of the prime-weighted factors \(p_k^{-\sigma}\)
as multiplication operators acting on \(\mathcal{H}_{\rm prime}\) and
derive basic operator-norm and Hilbert--Schmidt bounds.

\subsection{Prime-weighted diagonal operators}

\begin{definition}[Prime-weighted operator]
For a fixed exponent \(\sigma > 0\), define the diagonal operator
\(D_\sigma : \mathcal{H}_{\rm prime} \to \mathcal{H}_{\rm prime}\) by
\begin{equation}
  D_\sigma \ket{p_k} = p_k^{-\sigma} \ket{p_k},
  \qquad k\ge 1,
\end{equation}
and extended by linearity and continuity to all of
\(\mathcal{H}_{\rm prime}\).
\end{definition}

\begin{lemma}
The operator \(D_\sigma\) is bounded on \(\mathcal{H}_{\rm prime}\) for
every \(\sigma>0\), with operator norm
\(\|D_\sigma\|_{\mathcal{B}(\mathcal{H}_{\rm prime})}=1\).
\end{lemma}

\begin{proof}
For any \(\psi = \sum_{k} c_k \ket{p_k} \in \mathcal{H}_{\rm prime}\),
\begin{equation}
  \|D_\sigma \psi\|^2
  = \sum_{k=1}^{\infty} |p_k^{-\sigma} c_k|^2
  = \sum_{k=1}^\infty p_k^{-2\sigma} |c_k|^2.
\end{equation}
Since \(p_1 = 2\) is the smallest prime and \(p_k \ge 2\) for all \(k\),
we have \(0 < p_k^{-\sigma} \le 2^{-\sigma} < 1\). In particular,
\begin{equation}
  p_k^{-2\sigma} \le 1
  \quad\Rightarrow\quad
  \|D_\sigma \psi\|^2 \le \sum_{k=1}^{\infty} |c_k|^2 = \|\psi\|^2.
\end{equation}
Thus \(\|D_\sigma\|\le 1\). The norm is exactly \(1\) because we can
approach the bound by a unit vector supported on large primes, where the
weights \(p_k^{-\sigma}\) are arbitrarily close to \(0\), but the
supremum over all basis vectors is realized at \(k=1\) with
\(p_1^{-\sigma}<1\). By standard properties of diagonal operators with
bounded diagonal entries, we conclude \(\|D_\sigma\|=1\).
\end{proof}

\subsection{Hilbert--Schmidt bounds for prime-weighted kernels}

In various constructions it is convenient to interpret the prime
weights as defining an integral kernel on an auxiliary index space.
We collect the relevant bounds here.

\begin{definition}
Let \(\sigma>0\). We define the matrix kernel
\(K^{(\sigma)} = (K_{ij}^{(\sigma)})_{i,j\ge 1}\) on
\(\mathcal{H}_{\rm prime}\) by
\begin{equation}
  K_{ij}^{(\sigma)} \coloneqq p_i^{-\sigma}\,\delta_{ij}.
\end{equation}
The corresponding operator \(K^{(\sigma)}\) is simply \(D_\sigma\)
expressed in the prime basis.
\end{definition}

\begin{lemma}[Hilbert--Schmidt condition]
For any \(\sigma > \tfrac{1}{2}\), the operator \(K^{(\sigma)}\) is
Hilbert--Schmidt. In particular, its Hilbert--Schmidt norm satisfies
\begin{equation}
  \|K^{(\sigma)}\|_{\rm HS}^2
  = \sum_{i,j=1}^{\infty} |K_{ij}^{(\sigma)}|^2
  = \sum_{k=1}^{\infty} p_k^{-2\sigma} < \infty.
\end{equation}
\end{lemma}

\begin{proof}
By definition of the Hilbert--Schmidt norm,
\begin{equation}
  \|K^{(\sigma)}\|_{\rm HS}^2
  = \sum_{i,j=1}^\infty |p_i^{-\sigma}\,\delta_{ij}|^2
  = \sum_{i=1}^\infty p_i^{-2\sigma}.
\end{equation}
The series \(\sum p_i^{-s}\) over primes converges for any
\(s>1\), by comparison with the Riemann zeta function
\(\zeta(s) = \sum_{n\ge 1} n^{-s}\). In our case \(s=2\sigma\), so
\(s>1\) is equivalent to \(\sigma > \tfrac{1}{2}\). Therefore,
\(\sum p_i^{-2\sigma}<\infty\) whenever \(\sigma>\tfrac{1}{2}\), and
\(K^{(\sigma)}\) is Hilbert--Schmidt.
\end{proof}

\begin{corollary}
For \(\sigma>\tfrac{1}{2}\), the spectrum of \(K^{(\sigma)}\) consists
of a sequence of eigenvalues \(\{p_k^{-\sigma}\}_{k\ge 1}\) accumulating
only at \(0\), and \(K^{(\sigma)}\) is compact on \(\mathcal{H}_{\rm prime}\).
\end{corollary}

\begin{proof}
Diagonal operators with entries forming an \(\ell^2\)-sequence are
Hilbert--Schmidt and hence compact. The eigenvalues are the diagonal
entries \(p_k^{-\sigma}\), which converge to \(0\) as \(p_k\to\infty\).
\end{proof}

These bounds justify the choice \(\sigma>1/2\) in the main text for
embedding stability, and connect directly to the four-prime
normalization condition
\(\sum_{k=1}^{4} w_k p_k^{-\sigma} = 1\).

\section{Boundedness of the Prime-Structured Potential}

We next show that the explicit prime-structured potential used in the
tunneling problem defines a bounded multiplication operator on the
one-particle Hilbert space \(L^2(\mathbb{R})\), and that it is a
relatively bounded perturbation of the free Hamiltonian.

\subsection{Explicit potential and basic bounds}

Recall the explicit ansatz:
\begin{align}
  V_0(x) &= V_b\left(1 - \frac{x^2}{a^2}\right)^2 - E_0,
  \\
  V_{\rm prime}^{(4)}(x)
  &= \sum_{k=1}^{4} f_k(x)p_k^{-\sigma},
  \\
  f_k(x) &= A_k x^2 + B_k \cos(\omega_k x + \phi_k),
  \\
  \omega_k &= \frac{2\pi}{\log p_k},\qquad \phi_k = 0,
\end{align}
with parameters
\begin{equation}
  V_b=1,\quad a=2,\quad E_0=0.3,\quad
  A=[0.15,0.12,0.08,0.05],\quad
  B=[0.08,0.06,0.04,0.03],\quad
  \sigma=0.8.
\end{equation}
The total potential is
\begin{equation}
  V(x) = V_0(x) + V_{\rm prime}^{(4)}(x).
\end{equation}

\begin{lemma}[Growth bound]
There exists a constant \(C>0\) such that
\begin{equation}
  |V(x)| \le C(1 + x^4)
  \quad\text{for all } x\in\mathbb{R}.
\end{equation}
\end{lemma}

\begin{proof}
The background term \(V_0(x)\) satisfies
\begin{equation}
  |V_0(x)| = \left|(1 - x^2/a^2)^2 - E_0\right|
  \le C_0 (1 + x^4)
\end{equation}
for some constant \(C_0>0\), since the quartic dominates the polynomial
growth. For the prime-structured term,
\begin{align}
  |V_{\rm prime}^{(4)}(x)|
  &\le \sum_{k=1}^4 |f_k(x)|\,p_k^{-\sigma} \nonumber\\
  &\le \sum_{k=1}^4 \left(|A_k|\,|x|^2 + |B_k|\right)p_k^{-\sigma}
   \le C_1 (1 + x^2)
\end{align}
for some constant \(C_1>0\), using \(|\cos(\cdot)|\le 1\). Combining,
\begin{equation}
  |V(x)|
  \le C_0 (1 + x^4) + C_1 (1 + x^2)
  \le C(1 + x^4)
\end{equation}
for a suitable \(C\).
\end{proof}

\subsection{Self-adjointness and relative boundedness}

Let \(H_0\) be the free one-dimensional Schr\"odinger operator,
\begin{equation}
  H_0 = -\frac{1}{2m}\frac{d^2}{dx^2},
  \qquad \mathcal{D}(H_0) = H^2(\mathbb{R}),
\end{equation}
and define the full Hamiltonian
\begin{equation}
  H = H_0 + V(x)
\end{equation}
as an operator on \(L^2(\mathbb{R})\).

\begin{lemma}
The potential term \(V\) defines a relatively bounded perturbation of
\(H_0\) with relative bound zero. In particular, \(H\) is
self-adjoint on \(\mathcal{D}(H_0)\).
\end{lemma}

\begin{proof}
It suffices to show that multiplication by \(V(x)\) is relatively
bounded with respect to \(-d^2/dx^2\) with arbitrarily small relative
bound. By the established growth bound, \(V(x)\) is at most quartic in
\(|x|\). Standard results on Schr\"odinger operators with polynomial
potentials (see, e.g., \cite{reed_simon2}) imply that multiplication by
such a \(V\) is \(H_0\)-bounded with relative bound zero. Concretely,
for any \(\epsilon>0\) there exists \(C_\epsilon>0\) such that
\begin{equation}
  \|V\psi\| \le \epsilon \|H_0\psi\| + C_\epsilon \|\psi\|
  \quad\text{for all } \psi\in H^2(\mathbb{R}).
\end{equation}
By the Kato--Rellich theorem, \(H = H_0 + V\) is self-adjoint on
\(\mathcal{D}(H_0)\).
\end{proof}

\section{WKB Exponent and Functional Derivatives}

We now formalize the derivation of the functional derivative
\(\partial \ln T / \partial f_k\) used in the definition of the
effective topological response \(\Delta\chi_{\rm eff}\).

\subsection{Differentiation under the WKB integral}

For fixed mass \(m>0\) and energy \(E<V_{\max}\), define
\begin{equation}
  W(x) = V_0(x) - E > 0 \quad \text{on } [x_L,x_R],
\end{equation}
and let
\begin{equation}
  G(x) = \sum_{k=1}^4 f_k(x) p_k^{-\sigma}.
\end{equation}
The WKB exponent is
\begin{equation}
  S[f](E)
  \coloneqq \int_{x_L}^{x_R} \sqrt{2m\left(W(x)+G(x)\right)}\,dx,
\end{equation}
and the tunneling probability
\begin{equation}
  T(E) \approx \exp\left[-2 S[f](E)\right].
\end{equation}

We regard \(S[f]\) as a functional of the functions \(f_k(x)\).

\begin{proposition}
Assume \(W(x)\ge c>0\) on \([x_L,x_R]\) and \(f_k \in C^1([x_L,x_R])\).
Then the functional derivative of \(S[f]\) with respect to \(f_k\) is
\begin{equation}
  \frac{\delta S}{\delta f_k}(x;E)
  = \frac{m\,p_k^{-\sigma}}{\sqrt{2m\left(W(x)+G(x)\right)}}.
\end{equation}
Consequently, for small variations \(\delta f_k\),
\begin{equation}
  \delta S(E)
  = \sum_{k=1}^4 \int_{x_L}^{x_R}
    \frac{m\,p_k^{-\sigma}}{\sqrt{2m\left(W(x)+G(x)\right)}}\,
    \delta f_k(x)\,dx.
\end{equation}
\end{proposition}

\begin{proof}
We compute the Fr{\'e}chet derivative of \(S\). For a variation
\(\delta f_k(x)\), the induced variation \(\delta G(x)\) is
\(\delta G(x) = \delta f_k(x)p_k^{-\sigma}\). Then
\begin{align}
  \delta S(E)
  &= \int_{x_L}^{x_R}
     \frac{d}{d\epsilon}\bigg|_{\epsilon=0}
     \sqrt{ 2m\left(W(x)+G(x)+\epsilon\,\delta G(x)\right)}\,dx \nonumber\\
  &= \int_{x_L}^{x_R}
     \frac{1}{2}
     \frac{ 2m\,\delta G(x) }
          { \sqrt{2m\left(W(x)+G(x)\right)} }\,dx \nonumber\\
  &= \int_{x_L}^{x_R}
     \frac{m\,p_k^{-\sigma}}{\sqrt{2m\left(W(x)+G(x)\right)}}\,
     \delta f_k(x)\,dx.
\end{align}
Identifying the kernel of the linear functional in \(\delta f_k\) yields
the claimed expression for \(\delta S/\delta f_k\).
\end{proof}

\subsection{Derivative of \(\ln T\)}

Since \(T(E) \approx \exp[-2 S[f](E)]\), we find
\begin{equation}
  \ln T(E) \approx -2 S[f](E).
\end{equation}
Therefore,
\begin{equation}
  \frac{\delta \ln T}{\delta f_k}(x;E)
  = -2 \frac{\delta S}{\delta f_k}(x;E)
  = -2\,\frac{m\,p_k^{-\sigma}}{\sqrt{2m\left(W(x)+G(x)\right)}}.
\end{equation}
Integrating over \(x\) we recover the expression used in the main text
for the sensitivity of \(\ln T\) to changes in \(f_k\):
\begin{equation}
  \frac{\partial \ln T}{\partial f_k}
  = \int_{x_L}^{x_R}
    \frac{\delta \ln T}{\delta f_k}(x;E)\,dx
  = -2m p_k^{-\sigma}
    \int_{x_L}^{x_R}
    \frac{1}{\sqrt{2m\left(W(x)+G(x)\right)}}\,dx.
\end{equation}
In the regime where \(G\) is small compared to \(W\), this reduces to
the linearized form displayed in the main text.

\section{Four-Prime Closure Condition as a Normalization Constraint}

Finally we give a simple abstract formulation of the four-prime closure
condition
\(\sum_{k=1}^{4} w_k p_k^{-\sigma} = 1\) used to stabilize the
four-prime Newton holographic embedding.

\begin{definition}[Four-prime normalization functional]
Let \(\sigma>0\) and \(\mathcal{P}_4=\{p_1,p_2,p_3,p_4\}=\{2,3,5,7\}\).
Define the linear functional
\(\Lambda_\sigma : \mathbb{C}^4 \to \mathbb{C}\) by
\begin{equation}
  \Lambda_\sigma(\bm{w})
  = \sum_{k=1}^{4} w_k p_k^{-\sigma},
  \qquad \bm{w}=(w_1,\dots,w_4).
\end{equation}
\end{definition}

\begin{lemma}
For any \(\sigma>0\), \(\Lambda_\sigma\) is surjective and continuous
on \(\mathbb{C}^4\). In particular, the constraint
\(\Lambda_\sigma(\bm{w}) = 1\) defines a three-dimensional affine
subspace of \(\mathbb{C}^4\).
\end{lemma}

\begin{proof}
Continuity is immediate from the finite sum. Surjectivity follows by
choosing \(\bm{w}=(p_1^{\sigma},0,0,0)\), which gives
\(\Lambda_\sigma(\bm{w})=1\). Since the equation
\(\Lambda_\sigma(\bm{w})=1\) is a single nontrivial linear constraint in
a 4-dimensional vector space, the solution set is a codimension-one
affine subspace, hence 3-dimensional.
\end{proof}

In the context of the multiplicity bundle, \(\bm{w}\) encodes the
relative weights of the four prime modes. The normalization
\(\Lambda_\sigma(\bm{w})=1\) can be interpreted as fixing the effective
“prime charge” of the four-prime sector, ensuring a well-defined
embedding of the multiplicity field into the Newton holography
construction.

\bigskip

These results provide the operator-theoretic backbone for the
prime-structured tunneling module, ensuring that the constructions in
the main text rest on well-defined, bounded operators with controlled
spectral properties.


\nocite{*}
\bibliographystyle{unsrt}  
\bibliography{references}
\end{document}